El modelo de autorización del \dfn{backend} combina dos niveles: unos indicadores globales
de la cuenta, que viajan en el \acs{jwt}, y un rol con permisos en cada asignatura, que se
obtiene de la membresía. Entenderlo bien fue condición previa para diseñar la navegación,
porque el error más grave del prototipo previo era precisamente tratar el rol docente como
algo global (\cref{SEC:AN-BACKEND}).

\subsection{Nivel de la organización}

El servicio es \textit{multi-tenant}: cada usuario pertenece a una organización y las
consultas se filtran por ella. Dentro de ese marco, dos indicadores de la cuenta tienen
efecto en toda la aplicación:
\begin{description}
  \item[\texttt{is\_staff}.] Identifica al gestor de la organización. Gestiona usuarios,
        cursos académicos, asignaturas, configuración y registro de auditoría y, además,
        omite las comprobaciones de permisos de asignatura: puede actuar en cualquiera de
        ellas.
  \item[\texttt{is\_superadmin}.] Administración del sistema. Es el único perfil que puede
        crear organizaciones y el que permite consultas que atraviesan el filtro por
        organización.
\end{description}
Ninguno de los dos aparece en la respuesta del inicio de sesión; el \dfn{frontend} los lee
del \acs{jwt} al iniciar o renovar la sesión.

\subsection{Nivel de la asignatura}

Cada persona tiene, en cada asignatura en la que participa, uno de tres roles:
\texttt{COORDINATOR}, \texttt{TEACHER} o \texttt{STUDENT}. La misma persona puede coordinar
una asignatura y ser solo docente en otra, o ser estudiante en una tercera; por eso la
interfaz no puede tener una única ``vista de profesor'' y otra ``de estudiante'', sino que
debe decidir qué mostrar en cada asignatura y en cada examen según el rol que le devuelve
\texttt{GET /my-subjects/}.

Sobre el rol, el servicio define once permisos \texttt{can\_*} con un valor por defecto
(\cref{TAB:AN-PERMISSIONS}). La coordinación puede sobrescribirlos para una membresía
concreta (\texttt{PATCH /memberships/\{id\}/permissions/}), de modo que los permisos
efectivos pueden diferir de los de la tabla. El estudiante no tiene permisos de asignatura:
lo que puede ver de su propio examen lo deciden los seis permisos del examen
(\texttt{student\_permissions}: ver páginas, verlas tras la revisión, descargarlas, ver
anotaciones, \dfn{rubrica} y desglose) combinados con el estado del periodo de revisión.

\begin{table}[htbp]
  \centering
  \small
  \caption{Permisos de asignatura por defecto según el rol (\texttt{DEFAULT\_PERMISSIONS}
           del \dfn{backend}). El gestor (\texttt{is\_staff}) los tiene todos.}
  \label{TAB:AN-PERMISSIONS}
  \begin{tabular}{lcc}
    \hline
    \textbf{Permiso} & \textbf{COORDINATOR} & \textbf{TEACHER} \\
    \hline
    \texttt{can\_create\_exams}        & sí & no \\
    \texttt{can\_assign\_correctors}   & sí & no \\
    \texttt{can\_create\_rubric}       & sí & sí \\
    \texttt{can\_resolve\_issues}      & sí & no \\
    \texttt{can\_publish\_grades}      & sí & no \\
    \texttt{can\_manage\_reviews}      & sí & no \\
    \texttt{can\_view\_all\_instances} & sí & sí \\
    \texttt{can\_export\_grades}       & sí & no \\
    \texttt{can\_manage\_members}      & sí & no \\
    \texttt{can\_delegate\_perms}      & sí & no \\
    \texttt{can\_manage\_groups}       & sí & no \\
    \hline
  \end{tabular}
\end{table}

\subsection{Actores de la interfaz}

De la combinación de ambos niveles se derivan los cuatro actores para los que se diseñan
las pantallas:
\begin{description}
  \item[Gestor de la organización.] Cuenta con \texttt{is\_staff}. Administra la
        organización y puede intervenir en cualquier asignatura; también es quien realiza
        la ingesta manual de escaneos y atiende las páginas que no se han podido asociar a
        ningún examen.
  \item[Coordinación.] Rol \texttt{COORDINATOR} en una asignatura. Crea y configura
        exámenes, asigna correctores, resuelve incidencias de escaneo, publica notas,
        gestiona la revisión y administra miembros, grupos y permisos.
  \item[Docente.] Rol \texttt{TEACHER}. Corrige las \dfnpl{instancia} que tiene asignadas
        (o todas, con \texttt{can\_view\_all\_instances}), puede crear \dfnpl{rubrica} y
        resuelve las solicitudes de revisión sobre las \dfnpl{instancia} y problemas que
        tiene asignados por las reglas de corrección.
  \item[Estudiante.] Rol \texttt{STUDENT}. Consulta su nota cuando se publica y, durante
        el periodo de revisión y según los permisos del examen, su examen corregido, el
        desglose y la \dfn{rubrica}; puede solicitar la revisión de problemas concretos.
\end{description}
El superadministrador queda fuera del alcance de la interfaz: su única operación propia, la
creación de organizaciones, no tiene pantalla (\cref{SEC:AN-RF}).

\subsection{Consecuencias para el \textit{frontend}}

El \dfn{frontend} replica la tabla de permisos por defecto en
\texttt{src/lib/permissions.js} con un único fin: mostrar u ocultar acciones. La
autorización real sigue en el servicio, que vuelve a comprobar cada operación; ocultar un
botón no protege nada, solo evita ofrecer lo que va a fallar. Cuando los permisos efectivos
de una membresía difieren de los de la tabla, el servicio responde con un error de permiso
que la interfaz traduce y muestra. Que un usuario no pueda descubrir el identificador de su
propia membresía, y por tanto consultar sus permisos efectivos, es uno de los huecos del
contrato (\cref{SEC:AN-GAPS}).
